\section{Computer-assisted existence results}
\label{sec:certified-existence}

We distinguish the three computer-assisted existence results below from
the numerically observed ladder. The proof establishes a localized mode
at a parameter pair inside each specified box. It does not certify every
observed radiation minimum, an infinite ladder, or a nonlinear
nonradiating solution.

\subsection{Model and precise statement}
\label{sec:proof-model}
Consider the complex Klein--Gordon equation in three spatial dimensions,
\begin{equation}
 \partial_t^2\phi-\Delta\phi+U'(|\phi|^2)\phi=0,
 \qquad U(S)=S-S^2+\tfrac12 S^3.
 \label{eq:proof-kg}
\end{equation}
A stationary Q-ball has the form $\phi_0=e^{i\omega t}f(r)$, where
\begin{equation}
 f''+\frac{2}{r}f'=(1-\omega^2)f-2f^3+\tfrac32 f^5,
 \qquad f(0)=a,\quad f'(0)=0.
 \label{eq:proof-profile}
\end{equation}
We use a real orthonormal basis $Y_{\ell m}$ of spherical harmonics and
write a linear perturbation as
\begin{equation}
 \delta\phi=e^{i\omega t}
 \left[\frac{A(r)}{r}e^{i\rho t}
       +\frac{B(r)}{r}e^{-i\rho t}\right]Y_{\ell m}.
 \label{eq:proof-mode}
\end{equation}
In a complex harmonic basis the second sideband instead contains
$Y_{\ell m}^*$. With $S=f^2$, substitution in the linearization of
\eqref{eq:proof-kg} gives
\begin{equation}
 \begin{aligned}
 A''&=\left[1-4S+\tfrac92 S^2-(\omega+\rho)^2
                  +\frac{\ell(\ell+1)}{r^2}\right]A+CB,\\
 B''&=\left[1-4S+\tfrac92 S^2-(\omega-\rho)^2
                  +\frac{\ell(\ell+1)}{r^2}\right]B+CA,
 \qquad C=-2S+3S^2.
 \end{aligned}
 \label{eq:proof-radial-system}
\end{equation}
Regularity requires $A,B=O(r^{\ell+1})$. Define
\begin{equation}
 \kappa_0=\sqrt{1-\omega^2},\qquad
 k=\sqrt{(\omega+\rho)^2-1},\qquad
 \kappa_c=\sqrt{1-(\omega-\rho)^2}.
 \label{eq:proof-channel-numbers}
\end{equation}

\noindent\textbf{Computer-assisted existence statement.}
For each row of Table~\ref{tab:proof-cases}, the corresponding certificate
encloses a quadruple $(a_*,c_*,\rho_*,\omega_*)$ with $c_*>0$ such that
\eqref{eq:proof-profile} has a positive global solution satisfying
$r e^{\kappa_0r}f(r)\to c_*$, and
\eqref{eq:proof-radial-system} has a nonzero real solution regular at the
origin with
\begin{equation}
 |A(r)|+|B(r)|\leq K_J e^{-\kappa_c r}\quad(r\geq L),
 \qquad e^{\kappa_c r}B(r)\longrightarrow1.
 \label{eq:proof-decay}
\end{equation}
The resulting spatial fields in \eqref{eq:proof-mode} are smooth and
exponentially localized. The certification is
conditional on the correctness of the stated interval-arithmetic
implementation and its underlying arithmetic libraries.

\begin{table}[htbp]
\centering
\small
\begin{tabular}{ccll}
\hline
Case & $\ell$ & Enclosure of $\omega_*^2$ & Enclosure of $\rho_*$ \\
\hline
C0 & 0 & $0.797676787111865191750\pm4.8\,10^{-22}$
       & $1.744617544837398735781\pm6.8\,10^{-22}$ \\
C1 & 0 & $0.6851289044582160933\pm2.5\,10^{-20}$
       & $1.6903565973281434568\pm4.7\,10^{-20}$ \\
C2 & 1 & $0.7544960137826372331\pm7.8\,10^{-20}$
       & $1.826342030181069956\pm3.0\,10^{-19}$ \\
\hline
\end{tabular}
\caption{The three certified parameter boxes, displayed as outward-rounded
decimal enclosures. The proof uses the full exact dyadic boxes in
$(a,c,\rho,\omega)$, not these two-coordinate summaries. C0 is BEWEIS-1;
C1 and C2 are respectively T1 and T2 of BEWEIS-2. Their numerical ladder
assignments are $(\ell,n)=(0,1),(0,2),(1,1)$; the labels $n$ are not
separately certified nodal counts.}
\label{tab:proof-cases}
\end{table}

For the enclosed parameter pair and its fixed radial angular channel, every
solution of \eqref{eq:proof-radial-system} that is square integrable on
$[L,\infty)$ is a scalar multiple of the constructed solution. This is
\emph{geometric radial simplicity}. It neither establishes algebraic
simplicity of a dynamical generator nor excludes Jordan chains. For C2,
the three-dimensional $\ell=1$ angular multiplet remains degenerate.
The translation mode at $\rho=0$ does not belong to the fixed nonzero
frequency considered here.

The term \emph{embedded} is used operationally: $k^2>0$ and
$\kappa_c^2>0$, so a free asymptotic channel is open, while the constructed
mode is exponentially localized and has no oscillatory radiation tail.
We do not identify the essential spectrum of a fully specified operator
pencil in this proof. The perturbation is periodic after removing
$e^{i\omega t}$; the two laboratory frequencies are $\omega\pm\rho$,
with no claim about their rational commensurability.

\subsection{Matching to an infinite exterior}
\label{sec:proof-tails}
The matching radius $L$ is not a truncation at which the potential is
discarded. Both the profile and the mode are continued to infinity by
Volterra constructions. For the profile set
$\widetilde u(r)=r e^{\kappa_0r}f(r)$ and
$q_f=-2f^2+\tfrac32 f^4$. Its tail satisfies
\begin{equation}
 \widetilde u(r)=c+
 \int_r^\infty\frac{1-e^{-2\kappa_0(s-r)}}{2\kappa_0}
                q_f(s)\widetilde u(s)\,ds.
 \label{eq:proof-profile-volterra}
\end{equation}
For a chosen bound $\overline f>0$, put
\begin{equation}
 \begin{aligned}
 K_P&=\frac{(2+\tfrac32\overline f^2)e^{-2\kappa_0L}}
                  {4\kappa_0^2L^2},\quad
 \eta=2K_Pc^3,\quad m=c+\eta,\\
 \Lambda&=\frac{(6+\tfrac{15}{2}\overline f^2)m^2e^{-2\kappa_0L}}
                     {4\kappa_0^2L^2}.
 \end{aligned}
 \label{eq:proof-profile-constants}
\end{equation}
The conditions $m e^{-\kappa_0L}/L\leq\overline f$,
$K_Pm^3\leq\eta$, $\Lambda<1$, and $c>\eta$ give a positive unique
fixed point in $\|\widetilde u-c\|_\infty\leq\eta$. Common interval
majorants verify these inequalities over the entire parameter box.
They also bound the corrections to both profile matching data.

For the mode, write $dV=-4f^2+\tfrac92 f^4$ and use
$\widetilde A=e^{\kappa_cr}A$, $\widetilde B=e^{\kappa_cr}B$.
In the $\ell=0$ case the normalized Volterra kernels are
\begin{equation}
 K_A(r,s)=e^{-\kappa_c(s-r)}\frac{\sin(k(s-r))}{k},\qquad
 K_B(r,s)=\frac{1-e^{-2\kappa_c(s-r)}}{2\kappa_c}.
 \label{eq:proof-l0-kernels}
\end{equation}
The inhomogeneous free term is $(0,1)$. For $\ell=1$ use
\begin{equation}
 \begin{aligned}
 j(x)&=\sin(x)/x-\cos(x),\quad y(x)=-\cos(x)/x-\sin(x),\\
 d(u)&=e^{-u}(1+1/u),\quad g(u)=e^u(1-1/u),\\
 G_A(r,s)&=\frac{j(kr)y(ks)-y(kr)j(ks)}{k},\\
 G_B(r,s)&=\frac{d(\kappa_cr)g(\kappa_cs)-g(\kappa_cr)d(\kappa_cs)}{2\kappa_c}.
 \end{aligned}
 \label{eq:proof-l1-kernels}
\end{equation}
Here $K_{A,B}=e^{-\kappa_c(s-r)}G_{A,B}$ and the free normalized
term is $(0,1+1/(\kappa_cr))$. For $kL,\kappa_cL\geq1$, valid bounds are
\begin{equation}
 \begin{aligned}
 |K_A|&\leq a_1=\frac{1+(kL)^{-2}}{k},&
 |K_B|&\leq b_1=\frac{b_0}{2\kappa_c},\quad b_0=1+(\kappa_cL)^{-1},\\
 |e^{-\kappa_c(s-r)}\partial_rG_A|&\leq a_2=\sqrt{1+(kL)^{-2}},&
 |(\partial_r-\kappa_c)K_B|&\leq b_2=1+(\kappa_cL)^{-1}
                                      +\tfrac12(\kappa_cL)^{-2}.
 \end{aligned}
 \label{eq:proof-l1-majorants}
\end{equation}
In particular, the last bound concerns the normalized \emph{physical}
derivative, not $\partial_rK_B$ alone. Since
$|dV|+|C|\leq6f^2+\tfrac{15}{2}f^4$ is integrable, ordered Volterra
iteration converges with factorial majorants. With
$Q\geq\int_L^\infty(|dV|+|C|)\,ds$ and $\mu=\max(a_1,b_1)$, the
$\ell=1$ correction bounds are $(a_1,b_1,a_2,b_2)\varepsilon$, where
$\varepsilon=b_0(e^{\mu Q}-1)/\mu$.
Uniform convergence also supplies the parameter continuity needed below.

\subsection{Finite-dimensional existence certificate}
\label{sec:proof-krawczyk}
The four unknowns are $z=(a,c,\rho,\omega)$. Rigorous forward integration
computes the profile and two regular radial solutions $R_1,R_2$ from
the origin. For $\ell=1$ the initial variables are
$A=r^2\alpha$, $B=r^2\beta$; their regular equations have derivative
term $4\alpha'/r$ or $4\beta'/r$. For a source $F$, the origin
integral equation has weight
$r^2\int_0^1\tau(1-\tau^3)F(\tau r)\,d\tau/3$, of mass $r^2/10$.
This avoids starting a singular $2/r^2$ equation with approximate data.
Subsequent interval Taylor steps use checked complex a priori enclosures
and Cauchy remainder bounds; parameter variations enclose the required
Jacobian over the whole box.

The matching map is $H=H_0+E$. Its first two rows match profile value
and derivative. The last two match the Wronski forms
$W(\Psi,R_i)=\Psi\mathbin{\cdot}R_i'-\Psi'\mathbin{\cdot}R_i$,
using free Jost data in $H_0$ and bounded tail corrections in $E$.
For $\ell=0$, these two Wronskis equal the origin values of $\Psi$.
For $\ell=1$, regular solutions behave as $r^2e_i$ and normalized
singular solutions as $r^{-1}e_i/3$; their Wronski pairing is
$\delta_{ij}$. Vanishing of both Wronskis therefore removes both
singular coefficients.

Let $Z=z_0+[-\delta,\delta]$, let $D$ enclose $DH_0(Z)$, and suppose
$|E_i|\leq\epsilon_i$ on $Z$. A fixed dyadic preconditioner $Y$ is used
to verify
\begin{equation}
 z_0-YH_0(z_0)+(I-YD)[-\delta,\delta]
             +|Y|[-\epsilon,\epsilon]\subset\operatorname{int}Z,
 \qquad
 \max_i\sum_j |I-YD|_{ij}\frac{\delta_j}{\delta_i}<1.
 \label{eq:proof-krawczyk-inclusion}
\end{equation}
The continuous map $z\mapsto z-YH(z)$ maps $Z$ into itself, so Brouwer's
theorem gives a fixed point. The weighted matrix bound makes $Y$
invertible, hence the fixed point satisfies $H(z_*)=0$. Differentiability
of $E$ is not required, and uniqueness of this parameter quadruple is
not asserted. Profile positivity follows from positive interior
enclosures, a minimum-principle argument on the remaining finite
interval, and the positive exterior fixed point.

Finally, $2\kappa_0>\kappa_c$, verified for all three boxes, permits
construction of two bounded oscillatory comparison solutions at infinity.
Their mutual Wronskian is nonzero. Together with the decaying Jost
solution they span a three-dimensional subspace of the four-dimensional
radial Cauchy-data space. Every square-integrable solution is
Wronski-orthogonal to this subspace, whose symplectic complement is
exactly the decaying line. This proves the geometric simplicity stated
above.

\subsection{Scope of verification}
\label{sec:proof-verification-limits}
The defining certificates use $L=40$ for C0, $L=32$ for C1, and $L=36$
for C2, with 256-bit ball arithmetic. Additional runs use larger matching
radii and 320-bit arithmetic. C2's originally specified repeat at $L=40$
failed the inclusion test. A subsequent run with an enlarged $c$-radius
passed after reevaluating the bounds on the new box; this is an adaptive
additional certificate, not an originally successful repeat. The C2
existence conclusion rests on its first successful certificate alone.

Project-internal cross-house mathematical and code reviews found no blocking error in
the reviewed chains; the stated scope and harmonic conventions include
their corrections. Review and higher-precision reruns do not constitute
a second independent implementation. Correctness of Python--FLINT,
Arb/Acb, and the recorded execution remains part of the computational
trust base. For C0 the simplicity inequality was recorded and separately
checked, but was not included in the historical aggregate success flag;
the newer C1/C2 drivers include it explicitly. Source and input versions
are tied to the archived certificates by hashes.

The validated arithmetic uses Arb/Acb midpoint--radius enclosures
\cite{Johansson2017Arb}, accessed through Python--FLINT
\cite{PythonFlintDocs,FlintSoftware}.
The finite-dimensional test is a Krawczyk-type preconditioned inclusion
\cite{Krawczyk1969}; for the separately bounded continuous tail correction,
existence follows from the explicit Brouwer argument in
Section~\ref{sec:proof-krawczyk}, without asserting uniqueness of the
parameter quadruple. These references identify the methodological
background; the integration, tail bounds and inclusion hypotheses remain
the specific checks of the archived certificates.

These results establish three linear localized modes in specified
parameter boxes. They do not establish nonlinear radiation cancellation,
linear or nonlinear stability, algebraic spectral simplicity, the absence of other
modes or parameter zeros, or the existence of every member of the
empirical ladder.
